\section{Computational methods}
\label{sec:appendix}

\paragraph{Hardware and software.}
All computations were run on an Apple M4 Max (16 cores, 128\,GB) under macOS, with Python 3.9.6, NumPy 2.0.2, SciPy 1.13.1 (Quickhull only seeds the exact hull computation) and C code compiled with Apple clang 17.
Uniformly random arrangements were generated by Fisher--Yates shuffles (NumPy's default generator; \texttt{xoshiro256**} in C); the Monte Carlo of \Cref{tab:random_energy} uses $10^8$ arrangements per order.
All energies are evaluated from \eqref{eq:Efull}; the exact statements of \Cref{sec:theory} are checked by \texttt{verify\_theory.py} in exact integer or rational arithmetic (all $105{,}888$ single swaps of the $880+8$ squares, the $n=3$ enumeration, the witnesses), and the checks added in this revision (\Cref{cor:loly}, \Cref{cor:loly_null}, the witnesses, the data of \Cref{sec:lowmode_n4}, the product of inertia) by \texttt{check\_round1.py}.

\paragraph{Exact moments and variances.}
The means in \Cref{prop:random_energy_moments} are traces of the covariance matrix $\sigma_n^2G$.
The variances were computed as exact rationals for each $n\le45$ by expanding $\E[E^2]$ over the $15$ equality patterns of a cell quadruple, with the fourth-moment expectations of the uniform permutation expressed through the power sums $p_2,p_4$ of the centred values; the rational functions \eqref{eq:var_low} were then obtained by interpolation separately for odd and even $n$ and verified on $5$--$11$ further orders each, and against the exhaustive $n=3$ values.
Densities of the Gaussian-limit mixtures in \Cref{fig:random_energy} were computed by Imhof's inversion formula \cite{imhof1961}.

\paragraph{Meet-in-the-middle enumeration of $L_4$.}
The four statistics $T_1=\sum_{ij}j\,s_{ij}$, $T_2=\sum_{ij}i\,s_{ij}$, $T_3=\sum_i s_{ii}$, $T_4=\sum_i s_{i,n-1-i}$ are additive over cells.
The $16$ cells are split into the top two rows $A$ and the bottom two rows $B$.
For each of the $\binom{16}{8}=12{,}870$ value sets $V$ and each of the $8!$ bijections $V\to A$ the partial vector $(T_1,\dots,T_4)|_A$ is packed into a $27$-bit key; for the complementary set and each bijection $V^c\to B$ the key of (target $-$ partial) is formed; the two key lists are sorted and merged, and the sum over matching keys of the product of multiplicities is $|L_4|$.
This touches $2\times12{,}870\times40{,}320\approx1.04\times10^9$ partial vectors instead of $16!\approx2.1\times10^{13}$ arrangements and runs in $9$\,s on $16$ threads.
Matching pairs are expanded explicitly to obtain the $14{,}639{,}440$ members, which are then classified by row and column sums and by $D_4$-orbit.
Validation is described after \Cref{prop:L4}.

\paragraph{Exact hulls and their classification.}
Doubling all coordinates makes every point of a gem an integer vector.
Facet planes are recovered from an initial Quickhull triangulation as primitive integer equations, coplanar triangles are merged into faces, vertices are the points on at least three facet planes, edges are vertex pairs sharing two facets, and volumes are exact rationals; Euler's relation $V-E+F=2$ is checked for every gem.
Combinatorial types are the isomorphism classes of the bipartite vertex--facet incidence graphs, grouped by a Weisfeiler--Lehman hash and confirmed by explicit isomorphism tests.
Congruence is decided exactly by testing for a bijection of hull-vertex sets preserving all integer squared distances (grouping by the sorted distance multiset, then an edge-labelled graph isomorphism), with a floating-point Kabsch superposition as a cross-check.
For the Siamese proof (\Cref{prop:siamese-12}, \Cref{sec:app_siamese}) the finitely many polynomial identities and inequalities in $m$ are verified by exact integer polynomial arithmetic, an inequality $p(m)>0$ for all integers $m\ge2$ being certified by positivity of the coefficients of $p(2+t)$ or, failing that, by exact evaluation up to a Cauchy root bound.

\paragraph{Reproducibility.}
The scripts for every table and figure, the C enumerator and Monte Carlo code, the verification scripts, and the JSON files containing all reported numbers are provided as ancillary files, each directory with a summary of the numbers it produces and the commands that regenerate them.
The web application at \url{https://kylemath.github.io/MagicGemWebpage/} displays the $880$ essentially different order-four squares and renders the gem of a selected one.

\section{Variance closed forms}
\label{sec:app_variance}

For a uniformly random arrangement the variances of the two energies of \Cref{prop:random_energy_moments} are rational
functions of $n$ of the form
\begin{equation}
\begin{aligned}
\Var(\Elow)&=\frac{(n-1)(n^2+1)\,P_{\mathrm o}(n)}{129600\,n^{6}}\ (n\ \text{odd}),&\quad
&\frac{(n^2+1)\,P_{\mathrm e}(n)}{129600\,n^{4}}\ (n\ \text{even});\\
\Var(\Efull)&=\frac{(n-1)(n^2+1)\,Q_{\mathrm o}(n)}{180\,n^{6}}\ (n\ \text{odd}),&\quad
&\frac{(n-2)(n^2+1)\,Q_{\mathrm e}(n)}{180\,n^{5}}\ (n\ \text{even}),
\end{aligned}
\label{eq:var_low}
\end{equation}
with integer polynomials $P_{\mathrm o}$, $P_{\mathrm e}$ of degrees $11$, $10$ and leading coefficient $25$, and
$Q_{\mathrm o}$, $Q_{\mathrm e}$ of degrees $6$, $5$ and leading coefficient $5$, so that $\Var(\Elow)=n^8/5184+O(n^7)$ and
$\Var(\Efull)=n^3/36+O(n^2)$.  The polynomials are listed in the ancillary file \texttt{revision\_round2/appendix\_data.json}
(and in \texttt{random\_energy/variance\_fits.json}); they were obtained as described in \Cref{sec:appendix}, and at $n=3$ they
give $1532/2187$ and $398/2187$, the exhaustive values, at $n=4$ $14093/1440$ and $527/640$ (\Cref{tab:random_energy}).

\section{A witness for $k^*(5)=4$}
\label{sec:app_witnesses}

The arrangement
\[
\begin{pmatrix}1&22&11&24&7\\23&3&12&2&25\\18&6&19&17&5\\4&21&14&10&16\\20&9&15&8&13\end{pmatrix}
\]
has all row sums equal to $M(5)=65$ and column sums $66,61,71,61,66$, i.e.\ $\gamma=(1,-4,6,-4,1)=\bigl((-1)^j\binom4j\bigr)_j$;
by \eqref{eq:finite_diff} it satisfies the degree-$3$ conditions of \Cref{def:degree_k} without being semi-magic, so
$k^*(5)=4$.  The analogous witnesses for $6\le n\le13$ (\Cref{rem:kstar_search}) are in the ancillary file
\texttt{revision\_round2/appendix\_data.json}, together with the script that verifies them.

\section{Siamese gems: the remaining steps of the proof}
\label{sec:app_siamese}

We continue the proof of \Cref{prop:siamese-12} from the closed form $s(i,j)=n\langle i+j-m\rangle+\langle i+2j-2m\rangle+1$.

\emph{Step 2: affine reduction to five layers.}  In gem coordinates $x=j-m$, $y=m-i$, with $\langle t\rangle=t-n\lfloor t/n\rfloor$,
\[
z=\underbrace{n(x-y+m)+(2x-y+m)+1-\tfrac{n^2+1}{2}}_{L(x,y)\ \text{affine}}+\;n\,\zeta(x,y),\quad
\zeta:=-n\Big\lfloor\tfrac{x-y+m}{n}\Big\rfloor-\Big\lfloor\tfrac{2x-y+m}{n}\Big\rfloor .
\]
The map $\Phi(x,y,z)=(x,y,(z-L(x,y))/n)$ is an affine bijection, so $\Phi(\mathcal G(S))=\operatorname{conv}\{(x,y,\zeta)\}$ has the
same face lattice and vertex cells as $\mathcal G(S)$.  In the box $|x|,|y|\le m$ both floors lie in $\{-1,0,1\}$, the pairs
$(-1,1),(-1,0)$ and, by the central symmetry $\zeta(-x,-y)=-\zeta(x,y)$, $(1,-1),(1,0)$ are impossible, so $\zeta$ takes the
five values $\pm(n+1),\pm1,0$ on the lattice points of five convex lattice polygons (\Cref{fig:siamese-layers}a):
$T=\operatorname{conv}\{(-m,m),(-1,m),(-m,1)\}$ at $\zeta=n+1$ (cut out by $y\ge x+m+1$),
$U=\operatorname{conv}\{(-m,1-m),(-m,0),(-1,m-1)\}$ at $\zeta=1$ (by $2x+m+1\le y\le x+m$),
$M=\operatorname{conv}\{(-m,-m),(0,m),(m,m),(0,-m)\}$ at $\zeta=0$ (by $|y-2x|\le m$), and $D=-U$, $B=-T$ at $\zeta=-1,-(n+1)$.

\emph{Step 3: the twelve points and sixteen planes.}  Let
$v_1=(-m,m,n{+}1)$, $v_2=(-1,m,n{+}1)$, $v_3=(-m,1,n{+}1)$, $v_4=(-m,0,1)$,
$v_5=(-m,-m,0)$, $v_6=(0,-m,0)$, $v_7=(0,m,0)$, $v_8=(m,m,0)$, $v_9=(m,0,-1)$,
$v_{10}=-v_1$, $v_{11}=-v_3$, $v_{12}=-v_2$ (the cells listed in the statement), and consider the inequalities
\begin{align*}
&(1)\ \zeta\le 2m+2, \qquad (2)\ y\le m, \qquad (3)\ x\ge -m, \qquad (8)\ -2y+\zeta\le 2m,\\
&(4)\ -2(m{+}1)x+(2m{+}1)y-m\zeta\le m(2m{+}1),\\
&(5)\ 2(m{+}1)x-2(m{+}1)y+(2m{+}1)\zeta\le 2m(m{+}1),\\
&(6)\ 2m(m{+}2)x-(2m^2{+}2m{-}1)y+m(2m{+}1)\zeta\le m(2m^2{+}2m{-}1),\\
&(7)\ 2mx-y+m\zeta\le m(2m{-}1),
\end{align*}
together with their images $(1')$--$(8')$ under $(x,y,\zeta)\mapsto-(x,y,\zeta)$; let $Q$ be the polyhedron they define.
The script \texttt{siamese\_proof.py} verifies, as polynomial identities and inequalities in $m$ valid for all integers $m\ge2$
(\Cref{sec:appendix}): (a)~every vertex of $T,U,M,D,B$ at its height satisfies all $16$ inequalities, so
$\Phi(\mathcal G(S))\subseteq Q$; (b)~each plane contains exactly the following $v$'s, all others lying strictly inside:
\[
\begin{array}{llll}
(1)\colon v_1v_2v_3, & (2)\colon v_1v_2v_8v_7, & (3)\colon v_1v_3v_5v_4, & (4)\colon v_1v_4v_7,\\
(5)\colon v_2v_3v_6, & (6)\colon v_2v_6v_9, & (7)\colon v_2v_8v_9, & (8)\colon v_3v_5v_6,
\end{array}
\]
with centrally symmetric lists for $(1')$--$(8')$; (c)~the quadrilaterals are in convex position in the cyclic order given, so
the $16$ face polygons have $26$ distinct edges, each shared by exactly two polygons; (d)~each $v_k$ lies on three planes with
independent normals.  By (b) each plane supports $P:=\operatorname{conv}\{v_1,\dots,v_{12}\}$ in a facet with the stated vertices;
by (c) this family of facets is closed under edge-adjacency, and since the facet-adjacency graph of a $3$-polytope is connected
these are all the facets, so $P=Q$.  With (a), $P\subseteq\Phi(\mathcal G(S))\subseteq Q=P$, whose vertex set is
$\{v_1,\dots,v_{12}\}$ by (d); mapping back by $\Phi^{-1}$ gives the statement, with $(V,E,F)=(12,26,16)$ read off from (b)--(c).
For $m=1$ the triangles degenerate to points and the same argument gives the $8$-vertex Lo Shu gem.
